\chapter{Introduction}
\label{ch:introduction}

\section{Motivation}

Inference on phones, vehicles and embedded controllers is now commonly carried out by a
dedicated neural processing unit (NPU) rather than by the application CPU. Qualcomm's
Snapdragon platforms, for example, expose a Hexagon NPU through the Qualcomm AI Engine Direct
SDK, which provides a unified low-level API across the CPU, GPU and NPU of the chip
\cite{qualcomm_qnn}. To fit such hardware, a model is first compressed. The most common
compression is post-training quantization to 8-bit integers (INT8), which cuts weight storage
to a quarter of its 32-bit floating-point (FP32) size and allows integer-only arithmetic
\cite{jacob2018quantization,nagel2021whitepaper}. For large language models (LLMs), techniques
such as GPTQ \cite{frantar2023gptq} and SmoothQuant \cite{xiao2023smoothquant} make 8-bit and
lower-precision inference practical at billion-parameter scale, and model families such as
Llama~3.2, Qwen2.5 and Phi-3.5 include sizes intended for on-device use
\cite{llama3herd,qwen25report,phi3report}.

The deployed artefact is therefore a set of INT8 tensors held in DRAM, read by an accelerator
through direct memory access, and executed by a vendor runtime. That artefact has three
security properties worth protecting.

\begin{description}[leftmargin=*,style=nextline]
  \item[Integrity of the weights.] Rowhammer, first shown by Kim et al.\ \cite{kim2014flipping},
        lets unprivileged code flip bits in neighbouring DRAM rows. Hong et al.\ showed that most
        of the 19 image models they studied contain at least one parameter whose single bit flip
        drops accuracy by more than 90\% \cite{hong2019terminal}. Rakin et al.'s progressive
        bit-flip attack (PBFA) searches for the most damaging bits in quantized models
        \cite{rakin2019bfa}, and DeepHammer carried out such an attack end to end on real
        hardware, degrading accuracy to random guessing within minutes \cite{yao2020deephammer}.
        Model weights can also be modified before they reach the device, in the model supply
        chain. Feng and Tram\`er showed that a tampered pretrained model can expose the data
        on which it is later fine-tuned \cite{feng2024privacy}.
  \item[Integrity of the computation's inputs.] Adversarial examples, small perturbations of
        an input chosen to change the output, are a long-established threat
        \cite{goodfellow2015explaining,madry2018towards}. For LLMs the corresponding threats
        include prompt injection and attempts to bypass safety training.
  \item[Confidentiality of the weights and the training data.] Weights can be stolen through
        a prediction API \cite{tramer2016stealing} or read physically. Cold-boot attacks recover
        DRAM contents after power loss \cite{halderman2008lest}. Once an attacker can query a
        model or holds its weights, memorised training data becomes a target. Carlini et al.\
        extracted verbatim training sequences, including personal information, from GPT-2
        \cite{carlini2021extracting}. Nasr et al.\ scaled such extraction to production models
        \cite{nasr2023scalable}. Dai et al.\ reconstructed training text from activations
        exchanged in decentralised training \cite{dai2025activation}.
\end{description}

Defences exist for each of these. For weight integrity there are checksum and hash
signatures over weights \cite{li2021radar,javaheripi2021hashtag,guo2021modelshield}, honeypot
neurons \cite{liu2023neuropots} and weight reconstruction \cite{li2020defending}. For inputs
there are feature-space detectors such as the Mahalanobis-distance detector of Lee et al.\
\cite{lee2018simple}. For confidentiality there are trusted-execution-environment (TEE) designs
such as Amulet \cite{mao2025amulet}. These defences are typically designed and evaluated in
isolation, and most image-model defences were not designed with transformer LLM weight layouts
in mind.

\section{The problem studied}

This project asks whether a single, lightweight runtime layer can provide three things together
for a quantized model on an edge accelerator:

\begin{enumerate}[label=(\roman*),leftmargin=*]
  \item a \emph{cryptographic} guarantee that the weights being executed are the weights that
        were signed at release, with the modified tensor identified when they are not;
  \item a \emph{statistical} check that each inference's internal activations are consistent
        with clean behaviour, computed on INT8 values without dequantizing;
  \item a \emph{label and explanation} for what was detected, so that an operator knows
        whether to reload weights, drop an input or investigate a probe.
\end{enumerate}

The motivating observation is that the two checks fail in complementary ways. A hash check
sees nothing wrong with an adversarial input, because the weights are unchanged. An activation
monitor can miss a weight attack if the attacker also controls the input and uses it to keep
activations looking normal. This combined attack is called a \emph{synergy attack} in this
report. It is a construction used for evaluation, not an attack taken from the literature.

\section{Research questions}

The work is organised around five questions. Each is answered, within the limits of a
simulation-based study, in \cref{ch:evaluation}.

\begin{description}[leftmargin=*,style=nextline]
  \item[RQ1 — Weight integrity.] Does per-tensor SHA-256 hashing organised as a signed Merkle
        tree detect every simulated weight modification, including single-bit and
        least-significant-bit changes, and does it identify the modified tensor?
  \item[RQ2 — Activation anomalies.] Does a monitor calibrated on INT8 activations separate the
        simulated adversarial and extraction-probe inputs from clean ones, and at what
        false-positive rate?
  \item[RQ3 — Complementarity.] Does combining the two checks catch inputs that either check
        alone would miss, in particular synergy inputs?
  \item[RQ4 — Labelling.] How often does the rule-based classifier assign the intended threat
        label, and where does it confuse classes?
  \item[RQ5 — Cost.] What does each check cost per inference, and how does attestation cost
        scale to weight sets of language-model size?
\end{description}

\section{Contributions}

The contributions of this phase of the project are stated below. The report is careful to
separate what was built and measured from what is planned.

\begin{enumerate}[leftmargin=*]
  \item \textbf{A combined attestation-and-monitoring framework.} \textsc{EdgeGuard} is an
        open-source Python package of about 4{,}000 lines (\cref{ch:implementation}). It
        provides a signed SHA-256 Merkle tree over weight tensors, an INT8 activation monitor
        built on Welford's streaming statistics \cite{welford1962note}, and a rule-based
        classifier that produces per-layer explanations. Hash-based weight checking is not new
        in itself \cite{javaheripi2021hashtag,guo2021modelshield}; the contribution is the
        integration with activation monitoring and threat labelling, and its application to
        transformer weight layouts.
  \item \textbf{A simulation testbed.} The package includes simulators for bit-flip attacks
        under four strategies, supply-chain poisoning under three, adversarial inputs under
        three, training-data-extraction probes, physical weight extraction under three channel
        models, and synergy attacks. It also includes a generator of INT8 weight sets with the
        tensor shapes and HuggingFace naming of five open-weight language models. A suite of
        \NTests{} automated tests covers the package.
  \item \textbf{A measured evaluation with an explained false-positive rate.} All numbers in
        \cref{ch:evaluation} are generated by a single script from recorded runs. The clean
        false-positive rate of a fixed $4\sigma$ threshold is shown to follow from a
        multiple-comparisons argument (predicted \PredFPRZFour{}, measured \EmpFPRZFour{}). A
        Bonferroni-corrected threshold brings it to \EmpFPRBonf{} without losing any of the
        activation-level detections.
  \item \textbf{Evidence for complementarity, and a cost limit.} In simulation, the activation
        monitor alone flagged \MonSynergy{} of synergy inputs while the combined pipeline
        flagged \FlagSynergy{}. Full re-attestation of a 1.4\,GiB weight set costs about 1.6\,s
        on the test machine, which bounds where per-inference attestation is practical.
  \item \textbf{Engineering findings.} Two defects in the initial attack simulators were found
        and fixed during the evaluation. Both were integer-overflow errors in the handling of
        two's-complement sign bits, and both are of a kind that would silently distort a
        bit-flip study (\cref{sec:impl:bugs}).
\end{enumerate}

\section{Scope and limitations}

Three limitations should be stated at the outset, because they qualify every result.

First, \textbf{no real model is executed}. The weight sets are random INT8 tensors with
realistic shapes. The activations are drawn from distributions defined in the attack
simulators, not produced by running a network. Weight-level results do not depend on this,
because a hash detects a changed byte regardless of what the byte means. Activation-level
results do depend on it.

Second, \textbf{the detector and the attack generators were written together}. An activation
detector that achieves an AUC of 1.0 against attacks whose activation distributions were chosen
by the same authors demonstrates that the pipeline works end to end. It does not demonstrate
robustness against an adaptive attacker. \Cref{ch:discussion} discusses this in detail.

Third, \textbf{no NPU measurements are reported}. The QNN adapter in the package is a
duck-typed interface that has been exercised only with test doubles. Latencies are measured on
a two-core x86 cloud machine and are reported as such.

\section{Organisation of the report}

\Cref{ch:background} covers the background needed for the rest of the report: quantization,
the NPU software stack, transformer weight layouts, the cryptographic primitives used, and the
statistics behind the monitor. \Cref{ch:related} reviews prior work on bit-flip attacks and
defences, input anomaly detection, model and data extraction, and supply-chain attacks.
\Cref{ch:threat} states the threat model. \Cref{ch:design} presents the design and argues the
properties of each component. \Cref{ch:implementation} describes the implementation, the
attack simulators and the test suite. \Cref{ch:evaluation} reports the experiments.
\Cref{ch:discussion} interprets them and discusses limitations. \Cref{ch:conclusion} concludes
and sets out the plan for the next phase. The appendices give the configuration file, the API
reference, the full test inventory, per-seed results and reproduction instructions.
